\section{Reductions, complete problems, and constructive search}\label{sec:np}

The historical completeness constructions originate with Cook and Karp \cite{Cook1971,Karp1972}. Our 2SAT treatment uses the implication-graph criterion associated with Aspvall, Plass, and Tarjan \cite{AspvallPlassTarjan1979}, while proving only the explicit polynomial bound established below.

The results in this section sharpen the verifier and circuit framework into a
collection of explicit reductions and algorithms. Cook's original paper and
Karp's reduction paper are the historical sources for the completeness
framework; the proofs below give the finite constructions and witness maps
needed here. The implication-graph criterion for 2SAT is the one treated by
Aspvall, Plass, and Tarjan. We use a deliberately simpler polynomial algorithm,
not their sharper linear-time implementation. Bibliographic pointers are
listed in the source ledger accompanying this report.

\subsection{Conventions that affect correctness}

A Boolean instance has an explicitly specified finite variable universe
$0,\ldots,v-1$. Its representation lists those variables, or gives $v$ in
unary, so $v$ is at most the input length. An assignment is a bit vector of
length exactly $v$, including values of unused variables. This detail is
essential for counting. Ordinary syntactic formulas can enforce the same
convention by appending $x_i\vee\neg x_i$ for every designated unused variable.
A literal is a pair $(i,\epsilon)$, where $i<v$ and $\epsilon$ indicates whether
the bit is negated. A clause is a finite list of literals; an empty clause is
false and an empty conjunction is true. In $k$CNF every clause has width at
most $k$. Repeating literals pads nonempty short clauses to width exactly
three without changing assignments. Graphs below are finite, simple,
undirected graphs with numbered vertices and an explicit adjacency table.
Numbers in arithmetic decision problems are binary, and lists retain their
positions even if two entries have equal values.

A verifier presentation consists of a fixed deterministic machine $V$,
nonnegative-coefficient polynomial bounds $q,T$, and the decision bit
\[
 A_V(x)=\bigvee_{y<2^{q(n)}} V_T(x,y),\qquad n=|x|.
\]
Here $y$ is decoded as exactly $q(n)$ bits and $V_T$ is the acceptance bit after
$T(n)$ steps, with absorbing halting states. We can enlarge $T$ so that
$T(n)\ge n+q(n)+1$. There is no unbounded evaluation in this definition.
All instances are validated before processing. A reduction on valid instances
is extended to malformed input by a fixed answer convention and fixed valid
yes/no instances. In particular, the complement of a totalized language also
complements that convention; ``invalid'' is not silently identified with both
answers.

\begin{theorem}[N01: constructive algebra of reductions]\label{thm:N01}
Let $A,B,C$ be represented decision bits. If polynomial-time transducers $f,g$
satisfy $A(x)=B(f(x))$ and $B(z)=C(g(z))$, their composed transducer is a
polynomial-time reduction $A\leq_m^p C$. Identity and composition are explicit,
and polynomial-time decidability is preserved downwards. Logarithmic-space
many-one reductions have the same composition property.
\end{theorem}
\begin{proof}
Suppose $f$ uses at most $p(n)$ steps and outputs at most $p(n)$ bits, while
$g$ uses at most $r(m)$ steps on $m$ bits. A machine runs $f$, copies its output
into a new input area, and runs $g$. A multitape implementation costs at most
$c(p(n)+r(p(n))+n+1)^2$ for a fixed coding constant $c$; this deliberately
loose bound covers the copies and elementary simulation. Polynomial
composition and the displayed pointwise equalities prove the claim. To
combine a reduction with a decider, replace $g$ by that decider. No decision
about whether a machine is polynomial time is used: its bound is part of the
supplied presentation.

For logspace, storing $f(x)$ would be wrong. Instead, simulate $g$ with a
virtual input tape. Its input-head position uses $O(\log n)$ bits, because
$f$ has polynomial output length. Whenever $g$ asks for position $j$, restart
$f$ on $x$, count the emitted symbols to $j$, and provide that symbol. Save
$g$'s work tape while doing this. The sum of the two logarithmic work bounds
is logarithmic. A halting logspace transducer has polynomial time and output
length: ignoring its write-only output tape, it has only polynomially many
configurations, and a repeated configuration would repeat forever. One may
instead supply the resulting polynomial clock explicitly, which is the
preferred internal representation. Thus the recomputation also has a
polynomial bound. The invariant is equality of the simulated control,
work-tape, and virtual-input contents with the ordinary computation of $g$ on
$f(x)$, for all positions below the explicit clock. It is a bounded predicate.
\end{proof}

\begin{theorem}[N02: verifier and branching presentations of NP]\label{thm:N02}
Polynomial-time binary nondeterministic machines and polynomial-length,
polynomial-time verifiers define the same decision bits. The translations
send accepting witnesses to accepting witnesses by primitive-recursive maps.
\end{theorem}
\begin{proof}
For a binary branching machine with time bound $T(n)$, use a string of exactly
$T(n)$ bits to select the branch at each step. An absent transition rejects;
a deterministic transition ignores its bit; a halted computation ignores all
remaining bits. The deterministic verifier simulates these choices for the
supplied clock. An accepting path of length at most $T(n)$ supplies its choice
bits, padded by zeros. Conversely, the simulation supplies the selected
accepting path. The path code and the bit string have elementary bounds in
$T$, and equality of the two current configurations is the induction
invariant. A finite branching degree greater than two is compiled into a
fixed binary choice tree, rejecting unused leaves, at constant-factor cost.

For the other direction, a nondeterministic machine first guesses exactly
$q(n)$ witness bits and then runs $V$ for its supplied time. It accepts exactly
when some witness does. These are existential witness translations; they do
not assert a bijection between all padded paths and witnesses. Such a
bijection is needed only in the counting theorem below, where we fix the
verifier convention from the outset.
\end{proof}

\subsection{Circuit compilation and Cook--Levin}

A circuit is a list of input nodes followed by constant, NOT, binary AND, and
binary OR gates. Every non-input gate refers only to earlier nodes. A selected
node is its output. Evaluate in list order, recording one bit for each node.
After $i$ gates, the recorded prefix is the unique prefix satisfying those
$i$ gate equations for the fixed input bits: the next gate has already
computed arguments and a uniquely determined output. This bounded prefix
invariant proves both existence and uniqueness without an infinitary
valuation object.

\begin{theorem}[N03: uniform finite computation circuits]\label{thm:N03}
For a fixed verifier $V$ and a supplied polynomial clock $T(n)\ge n+q(n)+1$,
there is a polynomial-time constructor $x\mapsto C_x$ with $q(n)$ input bits
such that, for every $q(n)$-bit vector $y$,
\[
 C_x(y)=V_T(x,y).
\]
There is a constant $c_V$ with at most $c_V(T(n)+1)^2$ gates and a polynomial
bit-length bound for $C_x$. The construction is logspace for fixed $V,q,T$.
\end{theorem}
\begin{proof}
We describe the compiler sufficiently explicitly to remove an appeal to a
universal circuit theorem. Suppose $V$ has $k$ tapes, a fixed finite alphabet,
and a fixed finite control. During $T$ steps, every head stays within an
interval containing the initially occupied cells and $T$ additional cells
on each side. Give every tape an interval of length
$L=n+q(n)+2T+3$, with its origin shifted by $T+1$.

At time $t$ use one bit $S_{h,j,a}$ for each tape $h$, cell $j<L$, and alphabet
symbol $a$, one head-position bit $H_{h,j}$, and one bit $Q_r$ for each control
state. Legal encodings are one-hot in the symbol, head, and state positions.
The initial layer consists of constants and the input witness wires, or their
negations, so its one-hot encoding is explicit. For example a witness cell
has symbol-1 bit $y_i$, symbol-0 bit $\neg y_i$, and zero bits for the other
symbols. The fixed input $x$ is hardwired.

For each tape and symbol form the scanned-symbol bit
\[
 U_{h,a}=\bigvee_{j<L}(H_{h,j}\wedge S_{h,j,a}).
\]
Implement this OR by a binary chain. For each transition rule $\rho$, form a
selector $D_\rho$ by ANDing its state bit with its scanned-symbol bits. There
are only constantly many selectors, since $V$ is fixed. The next-state bit is
the OR of selectors entering that state. For a cell-symbol bit use
\[
 S'_{h,j,a}=(\neg H_{h,j}\wedge S_{h,j,a})
 \vee\left(H_{h,j}\wedge
       \bigvee_{\rho\text{ writes }a\text{ on }h}D_\rho\right).
\]
If rule $\rho$ moves tape $h$ by $d_{\rho,h}\in\{-1,0,1\}$, define
\[
 H'_{h,j}=\bigvee_\rho(D_\rho\wedge H_{h,j-d_{\rho,h}}),
\]
with out-of-interval head bits set to zero. The reserved boundary cells ensure
that an actually reached head is never lost during the $T$ layers. Absorbing
halting states are implemented by rules that retain the tape and head and
repeat the state.

Every displayed conjunction or disjunction has constant length except the
$L$-term scanned-symbol chains. Consequently one layer uses at most $c_V L$
gates. Repeat the layer $T$ times and output the accepting-state bit. The
invariant says that every state, head, and cell-symbol bit in layer $t$ is the
corresponding bit of the $t$-step machine configuration, and that all one-hot
conditions hold. The initial layer proves the base case. A one-hot encoding
makes exactly one rule selector true; the three update formulas then give
precisely the tape, head, and state update. Bounded induction on $t\le T$
proves the claim for every witness, the latter quantified below $2^{q(n)}$.

Since $L\le 3T+3$, the gate bound is quadratic. Each node address has
$O(\log(T+2))$ bits, so the encoded circuit length is
$O_V((T+1)^2\log(T+2))$. A simple implementation that repeatedly scans its
partial output still has, for example, a bound $c'_V(T+2)^6$; no unit-cost
large-integer assumption is needed. More efficiently, preallocate a fixed
number of wire slots for each time, tape, cell, and local template. Unused
slots are constant gates. Loop counters and address arithmetic use
$O(\log(T+2))=O(\log(n+2))$ bits, and the bits of $x$ can be reread by an
input scan. This emits the circuit in logspace. The circuit's large values
are gate identifiers of logarithmic bit length, not evaluations of arbitrary
primitive-recursive syntax.
\end{proof}

\begin{theorem}[N04: Tseitin constraints have unique extensions]\label{thm:N04}
There is a linear-size conversion of an acyclic circuit $C$ to a 3CNF formula
$F_C$, with the original inputs retained as designated variables, such that
restriction and gate evaluation give inverse maps between satisfying
assignments of $F_C$ and inputs on which $C$ outputs one.
\end{theorem}
\begin{proof}
Use one variable for every circuit node. For a gate with output $z$, insert
the following clauses:
\[
\begin{array}{c|l}
 z=a\wedge b & (\neg z\vee a),\ (\neg z\vee b),\ (z\vee\neg a\vee\neg b)\\
 z=a\vee b & (z\vee\neg a),\ (z\vee\neg b),\ (\neg z\vee a\vee b)\\
 z=\neg a & (z\vee a),\ (\neg z\vee\neg a)\\
 z=1 & (z)\\
 z=0 & (\neg z).
\end{array}
\]
Append a unit clause asserting the output variable. For the AND line, if
$z=1$, the first two clauses force $a=b=1$; if $z=0$, the last clause says
that at least one input is zero. Conversely, those two cases verify all
three clauses whenever $z=a\wedge b$. The OR proof is dual; the NOT and
constant lines are two-case or one-case checks. These are finite bit case
splits, hence intuitionistically valid.

Gate evaluation extends every accepting input to an assignment satisfying
all clauses. Conversely, induction along the topological gate list shows
that a satisfying assignment has exactly the evaluation bit at every node.
Its output unit clause therefore forces acceptance. That induction also
proves uniqueness of the extension, including gates not connected to the
output. There are at most $3G+1$ clauses for $G$ gates. If the encoding derives
its variable universe solely from literal occurrence, add one tautology per
input; this changes the bound to $3G+q+1$ and prevents loss of unused inputs.
No new choices are introduced by repeating literals to reach width three.
All assignment maps are finite array operations, bounded by $2^{q+G}$.
\end{proof}

\begin{theorem}[N05: SAT, circuit SAT, and 3SAT are NP-complete]\label{thm:N05}
Every polynomial verifier has an explicit polynomial many-one reduction to
3SAT, with primitive-recursive forward and reverse witness maps. SAT and
circuit SAT are NP-complete as well.
\end{theorem}
\begin{proof}
Apply N03 to $V,x$ and N04 to $C_x$. An accepting witness $y$ gives the unique
satisfying assignment by gate evaluation. A satisfying assignment restricts
to $y$ and the simulation invariant gives $V_T(x,y)=1$. Thus
$A_V(x)=\operatorname{SAT}(F_{C_x})$ as equality of Boolean decisions.
The two compilers have polynomial composition. Membership in NP is witnessed
by a bit vector assigning every designated variable: evaluating a formula,
circuit, or CNF is polynomial in its explicit encoding. General SAT receives
the 3CNF formula unchanged, while the circuit target uses $C_x$ directly.
Every claim concerns a supplied verifier presentation; there is no search for
an unknown polynomial bound.
\end{proof}

\begin{theorem}[N06: circuit value is P-complete under logspace reductions]\label{thm:N06}
The problem of evaluating an explicitly supplied acyclic Boolean circuit with
specified input bits is in P and is hard for P under logarithmic-space
many-one reductions.
\end{theorem}
\begin{proof}
Evaluate the nodes in order, recording their bits. Array access by scanning
gives a polynomial tape-time bound, so no random-access assumption is needed.
For a fixed polynomial-time decider $M$ and its clock, apply N03 with no
witness inputs. The resulting constant-input circuit has value $M(x)$. The
last paragraph of N03 gives a logspace constructor. Its correctness is the
same layer invariant, so every presented P language reduces to circuit value.
\end{proof}

\subsection{Graph and arithmetic complete problems}

\begin{theorem}[N07: CLIQUE is NP-complete]\label{thm:N07}
Given a 3CNF $F$ with $m$ clauses, one can construct in polynomial time a graph
$G_F$ such that $F$ is satisfiable exactly when $G_F$ has an $m$-vertex clique.
\end{theorem}
\begin{proof}
Make a vertex for each literal occurrence, remembering both its clause and its
literal. Two vertices are adjacent exactly when they come from different
clauses and their literals are not complements. There are at most $3m$
vertices; a double scan writes the adjacency matrix in polynomial time.
From a satisfying assignment, select the first true literal occurrence in
each clause. Selected vertices are in different clauses and cannot carry
complementary literals, since both are true. They therefore form a clique.

Conversely, a clique contains at most one occurrence from each clause. A
clique of size $m$ therefore contains exactly one from each of the $m$
clauses. Assign a variable the value requested by its selected occurrence,
if there is one, and zero otherwise. Pairwise compatibility ensures that two
selected occurrences never request different values. Every clause then has
a true selected literal. Selection, compatibility checking, and assignment
formation are bounded scans. An empty clause makes an $m$-clique impossible;
if $m=0$, the empty assignment condition and empty clique both satisfy their
respective instances. Membership of CLIQUE in NP follows by checking a list
of $k$ distinct vertices and its at most $k^2$ pairs. Values $k$ larger than
the vertex count are rejected immediately.
\end{proof}

\begin{theorem}[N08: independent set and vertex cover are NP-complete]\label{thm:N08}
For a graph $G$ on $N$ vertices and $0\le k\le N$, the following equivalent
witness transformations are polynomial:
\[
\begin{split}
G\text{ has a }k\text{-clique}
 &\ \Longleftrightarrow\ \overline G\text{ has an independent set of size }k,\\
G\text{ has an independent set of size }k
 &\ \Longleftrightarrow\ G\text{ has a vertex cover of size }N-k.
\end{split}
\]
Consequently INDEPENDENT SET and VERTEX COVER are NP-complete.
\end{theorem}
\begin{proof}
For distinct vertices, complementing the adjacency bit changes ``all selected
pairs are edges'' into ``all selected pairs are nonedges.'' This leaves the
selected vertex list unchanged. For the second equivalence, let $S$ be the
selected independent set and output its vertex complement $C$. An edge
uncovered by $C$ would have both endpoints in $S$, contradicting independence.
Conversely, two adjacent vertices in the complement of a cover would give an
uncovered edge. Cardinalities satisfy $|C|=N-|S|$ by a bounded sum of membership
bits. For the common ``at least $k$'' and ``at most $N-k$'' conventions, trim
an oversized independent set by taking its first $k$ members; enlarge a
smaller cover when necessary. The input bounds make these operations finite.
Both target problems have polynomial verifiers checking their membership
bits and edges. Compose these transformations with N07.
\end{proof}

\begin{theorem}[N09: a parsimonious 3SAT-to-SUBSET-SUM reduction]\label{thm:N09}
For every 3CNF on $v$ designated variables with $m$ clauses, there is a
polynomial-time binary SUBSET-SUM instance with $2v+2m$ indexed positive
integers. Satisfying assignments and subsets summing to the target are in
explicit bijection. In particular, SUBSET SUM is NP-complete.
\end{theorem}
\begin{proof}
Use base eight, with $v$ variable columns and $m$ clause columns. For each
variable $x_i$, create two indexed numbers $a_i^0,a_i^1$. Both have digit one
in variable column $i$ and zero in every other variable column. In clause
column $j$, the digit of $a_i^b$ is the number of occurrences in clause $j$
that are made true by assigning $x_i=b$. This digit is at most three; repeated
literals are counted with their multiplicity. For each clause $j$, create
slack numbers $s_j^1,s_j^2$ with sole nonzero digit respectively one and two
in clause column $j$. The target has digit one in every variable column and
digit four in every clause column.

In any subset of the constructed list, a variable-column sum is at most two.
The total contribution from all variable numbers to a clause column is at
most three, because a literal occurrence contributes to exactly one of the
two numbers for its variable. The two slack numbers contribute at most three
more. Thus every column sum is at most six, strictly below eight. No carries
can occur in any subset sum. Equality with the target can consequently be
checked column by column, a statement proved by repeated quotient and
remainder by eight.

The variable columns force selection of exactly one of $a_i^0,a_i^1$ for
each $i$. These choices specify an assignment. If $t_j$ is the number of true
literal occurrences in clause $j$, then $0\le t_j\le3$, and its target column
requires a slack contribution $4-t_j$. The available slack sums are
$0,1,2,3$, each produced by a unique subset of the one- and two-slack numbers.
Hence the column can be completed exactly when $t_j\ge1$, and the completion
is unique. For $t_j=1,2,3$, select slack of total $3,2,1$, respectively.
This gives both directions and the bijection. If a clause is empty, its
column can never reach four, exactly as required.

Each integer has at most $3(v+m)$ binary bits because base eight is $2^3$.
The whole list has polynomial bit length and is written by nested bounded
loops over variables, clauses, and occurrences. An indexed subset is a bit
vector of length $2v+2m$; binary addition and comparison verify its sum in
polynomial time. Equal-valued entries remain distinct positions, so the
bijection also holds for unused variables and duplicate clauses. For the
zero-column case the list and target are empty and zero; there is one empty
subset and one zero-variable assignment. All nondegenerate generated numbers
are positive.
\end{proof}

\begin{theorem}[N10: PARTITION and 0--1 KNAPSACK are NP-complete]\label{thm:N10}
SUBSET SUM has explicit polynomial reductions to PARTITION and to the
0--1 knapsack decision problem with nonnegative integer weights and profits.
\end{theorem}
\begin{proof}
Let the subset-sum input be indexed numbers $a_1,\ldots,a_r$, target $T$, and
$S=\sum_i a_i$. If $T>S$, output a fixed no-instance. If $S=0$, decide the
instance directly and output a fixed yes- or no-instance. Otherwise
$S>0$ and $0\le T\le S$. Append the numbers $2S-T$ and $S+T$. The total is
$4S$, so equal partition asks for a sublist of sum $2S$. A subset of the
original list of sum $T$, together with $2S-T$, gives such a partition.
Conversely, a sublist summing to $2S$ contains exactly one appended number:
without them its sum is at most $S$, and with both its sum is at least $3S$.
If it contains $2S-T$, its original entries sum to $T$. If it contains $S+T$,
its original entries sum to $S-T$, so the complementary original entries sum
to $T$. The extra numbers have polynomial bit length.

For knapsack, give item $i$ weight and profit both $a_i$, capacity $T$, and
required profit $T$. The inequalities $\sum a_i\le T$ and $\sum a_i\ge T$
force equality. Conversely a subset sum of $T$ meets both constraints.
Checking subset sums, weights, and profits uses polynomial binary arithmetic,
so both target problems lie in NP.
\end{proof}

\begin{theorem}[N11: 0--1 integer feasibility is NP-complete]\label{thm:N11}
Feasibility of a finite system of integer linear inequalities with every
variable constrained to $\{0,1\}$ is NP-complete, even for the inequalities
arising from clauses of width at most three.
\end{theorem}
\begin{proof}
Replace each positive literal $x_i$ by the integer $x_i$ and each negative
literal by $1-x_i$. For each clause require the sum of its literal terms to
be at least one. On a bit assignment this sum counts exactly its true literal
occurrences. Therefore the same assignment satisfies the formula if and only
if it satisfies the system. The expansion has polynomial size, with bounded
coefficients even when a clause repeats a variable. An empty clause becomes
$0\ge1$. Membership in NP follows by guessing the bit vector and evaluating
all integer sums. If coefficients have $b$ bits and there are $r$ terms,
intermediate sums have at most $b+\lceil\log_2(r+1)\rceil+2$ bits, which
supplies the binary-arithmetic bound. This statement does not silently assert
a certificate bound for unrestricted integer variables.
\end{proof}

\begin{theorem}[N12: elementary NP presentations of neighboring problems]\label{thm:N12}
Graph isomorphism, Hamiltonian path, Hamiltonian cycle, and the bounded-factor
problem have explicit polynomial verifiers. The latter asks, on $N\ge2$,
whether $N$ has a divisor $d$ with $2\le d\le b$ and $d<N$.
\end{theorem}
\begin{proof}
For graph isomorphism on $r$ vertices, a witness lists $r$ vertex images.
Check that every image is below $r$, no two are equal, and all $r^2$
adjacency bits are preserved. The first two properties make the map a
permutation: a missing image would leave $r$ distinct images in $r-1$ slots,
contradicting the finite counting lemma. The witness uses
$O(r\log(r+2))$ bits and all checks are polynomial. Graphs of different sizes
are immediately rejected.

A Hamiltonian-path witness is a permutation of the vertices with an edge
between consecutive entries. For a cycle also check the last-to-first edge,
using the chosen convention for graphs of fewer than three vertices. These
verifiers directly express the defining objects. No NP-hardness of either
Hamiltonian problem is inferred merely from membership.

For bounded factor, guess the binary integer $d$, check the inequalities, and
compute $N\bmod d$ by binary long division. Its bit length is bounded by that
of $N$, and repeated shift, comparison, and subtraction give polynomial time.
The condition is equivalent to having a prime factor in the same range:
if a divisor $d\ge2$ is supplied, boundedly select its least divisor
$e\ge2$. A nontrivial factorization of $e$ would give a smaller divisor of
$d$, so $e$ is prime; also $e\le d$ and $e\mid N$. This equivalence does not
require computing $e$ in polynomial time, since the NP witness remains $d$.
It proves membership only, not a polynomial factoring algorithm or a
coNP presentation.
\end{proof}

\subsection{Tractable satisfiability fragments}

\begin{theorem}[N13: constructive 2SAT criterion and polynomial algorithm]\label{thm:N13}
A 2CNF formula without empty clauses is satisfiable if and only if no variable
and its complement are mutually reachable in its implication graph. There is
a polynomial algorithm returning either a satisfying assignment or a
variable together with the two contradictory implication paths.
\end{theorem}
\begin{proof}
Replace a unit clause $a$ by $a\vee a$. For each clause $a\vee b$, put edges
$\neg a\to b$ and $\neg b\to a$ on the $2v$ literal vertices. Every edge
$u\to w$ is paired with $\neg w\to\neg u$. An assignment satisfies the
formula exactly when it respects all implications. Such an assignment
respects every path by induction on its length. If $x_i$ and $\neg x_i$
are mutually reachable, either value of $x_i$ then forces its complement to
have that same truth value. Splitting on the bit of $x_i$ proves failure.

For the converse and an actual algorithm, compute reflexive reachability by
the recurrence
\[
 R_0(u,w)=[u=w]\vee E(u,w),\qquad
 R_{j+1}(u,w)=R_j(u,w)\vee(R_j(u,j)\wedge R_j(j,w)).
\]
The invariant says that $R_j$ describes paths whose intermediate vertices
lie below $j$. Path concatenation proves the forward update. For the reverse
update, erase repeated vertices from a path first; a simple path uses $j$
at most once as an intermediate vertex, and splitting there proves the
recurrence. Length is always bounded by $2v$. Recording a predecessor or
split index at the first successful update reconstructs paths by bounded
loops. Thus the table and any path certificates have explicit elementary
bounds; there is no unbounded reachability search.

Define $u\sim w$ when both $R(u,w)$ and $R(w,u)$ hold. Reflexivity,
symmetry, and transitivity follow from path operations. Represent each class
$C$ by its least literal number $r(C)$. Complementing every literal and
reversing every edge carries a path $u\leadsto w$ to
$\neg w\leadsto\neg u$; consequently complementation induces an involution
$C\mapsto\bar C$ on the classes. By hypothesis $C\ne\bar C$.

Here is an explicit ordering device that avoids an unproved topological-sort
subroutine. Let $d(C)$ count the classes $D$ with a path $D\leadsto C$,
including $C$ itself. Put
\[
 B=2v+1,\qquad K(C)=B\,d(C)+r(C).
\]
Different classes have different keys, since their representatives are below
$B$. If $C\leadsto D$ and $C\ne D$, every predecessor of $C$ is a predecessor
of $D$, while $D$ is not a predecessor of $C$ (otherwise the classes coincide).
Thus $d(D)\ge d(C)+1$, and $K(D)>K(C)$. This is a finite injection and
strict-cardinality argument over the list of representatives.

Assign a literal $u$ the value one exactly when
$K([u])>K([\neg u])$. The unequal keys make complementary literals receive
opposite bits. An edge $u\to w$ gives
$K([u])\le K([w])$ and, from its paired edge,
$K([\neg w])\le K([\neg u])$. If $u$ were true and $w$ false, these
inequalities would yield
\[
 K([\neg u])<K([u])\le K([w])
 <K([\neg w])\le K([\neg u]),
\]
an impossibility. Every implication is therefore respected, and the
assignment satisfies the formula. All quantities come from the finite
reachability table. Floyd's recurrence uses $O(v^3)$ Boolean operations;
representative scans and predecessor counts are polynomial too. Even a
sequential tape implementation with repeated table scans has a polynomial
bound. If there is an empty clause, return its direct inconsistency
certificate before constructing the graph. The case $v=0$ is handled by the
empty conjunction convention.
\end{proof}

\begin{theorem}[N14: Horn satisfiability has a least-model algorithm]\label{thm:N14}
Satisfiability of CNF formulas with at most one positive literal per clause is
polynomial-time decidable. On satisfiable input an explicit algorithm returns
the coordinatewise least satisfying assignment.
\end{theorem}
\begin{proof}
Write each clause as a rule $x_{i_1}\wedge\cdots\wedge x_{i_r}\to h$,
where $h$ is either a variable or false. Repeated antecedents do no harm.
Start with all variables false. In a fixed order scan all rules whose heads
are variables; whenever every antecedent is currently true, set the head
true. Repeat the complete scan $v$ times. A change adds a previously absent
true variable, and no variable is removed. If a full scan makes no change,
all later scans are identical. If the state after $v$ scans were not closed,
each of those scans would have made a change and a further change would be
possible, requiring more than $v$ distinct additions. Finite counting rules
this out. Thus the final state is closed under all positive-head rules.

At every addition, every satisfying assignment must make the newly added
variable true: it already makes all antecedents true by the induction
hypothesis and must respect this rule. Hence the computed true set is a
subset of every model. If any false-head rule now has all antecedents true,
every model would violate that rule, so return unsatisfiable. Otherwise all
rules are satisfied by the computed assignment; it is therefore a model,
and its containment in every model proves leastness. Facts are rules with
empty antecedent; an empty clause is the false-head rule with empty
antecedent and is rejected. At most $v$ scans of the input and polynomial
array accesses are needed. The invariant quantifies over assignments below
$2^v$ only; this makes its minimality proof an instance of bounded induction.
\end{proof}

\subsection{Search, conditional collapse, and counting}

\begin{theorem}[N15: constructive SAT search-to-decision]\label{thm:N15}
A supplied polynomial-time SAT decider yields a polynomial-time algorithm
which returns a satisfying assignment for any satisfiable formula and returns
``none'' for an unsatisfiable one. The same supplied decider yields such a
search algorithm for every polynomial verifier presentation.
\end{theorem}
\begin{proof}
Run the decider once on $F$. If it says no, return none. Otherwise maintain a
prefix assignment $a_0,\ldots,a_{i-1}$ for which the restricted formula is
satisfiable. Query whether this restricted formula with $x_i=0$ is
satisfiable. If yes, append zero. If no, append one. To justify the second
case constructively, take a witness for the currently true finite SAT bit;
its $i$th bit is either zero or one. Zero contradicts the no-answer, so one
supplies a witness for the selected restriction. The invariant is thus
preserved. After $v$ steps all variables are assigned and the restricted
formula's true value means exactly that this assignment satisfies $F$.

There are $v+1$ calls and each restricted encoding has length at most a fixed
polynomial in the original length. One may retain the original formula and
append unit clauses to specify the prefix. If the decider runs in time $p$,
a bound such as $(v+1)p(c(|F|+v+1)^2)+c'(|F|+v+1)^4$ suffices. The algorithm
never searches exponentially for the witness used in the correctness proof;
it uses the decider and finite bit case analysis.

For a verifier $V$, first construct $F_{C_x}$ by N03--N04, find its satisfying
assignment, and restrict to its designated witness-input variables. The
simulation theorem proves that this is an accepting witness. An initial
no-answer correctly certifies that none exists. The resulting algorithm is
a total search procedure with a distinguished failure result, not a claim
that every arbitrary multivalued selection becomes a prescribed function.
\end{proof}

\begin{theorem}[N16: the constructive completeness implication]\label{thm:N16}
A polynomial SAT decider implies a polynomial decider and a polynomial witness
finder for every presented NP problem. Conversely, if every NP presentation
has a polynomial decider, SAT does. Thus the standard statement
$\mathrm P=\mathrm{NP}$ is equivalent to $\mathrm{SAT}\in\mathrm P$ under the
numerical class-presentation convention of this report.
\end{theorem}
\begin{proof}
The forward decision transformation composes the supplied SAT machine with
the two compilers of N05. The witness-finding transformation is N15. These
are explicit functions of the SAT descriptor and the verifier descriptor;
their polynomial bounds are obtained by composition of the supplied
polynomials. The reverse implication specializes the class-equality
hypothesis to the concrete polynomial SAT verifier. The known inclusion
$\mathrm P\subseteq\mathrm{NP}$ supplies the other half of the equality.
These implications neither decide SAT in polynomial time nor settle the
class equality.
\end{proof}

\begin{theorem}[N17: an NP-complete coNP problem collapses NP and coNP]\label{thm:N17}
If a supplied NP-complete problem $K$ also has a coNP presentation, then
$\mathrm{NP}=\mathrm{coNP}$. In particular, $\mathrm P=\mathrm{NP}$ implies
$\mathrm{NP}=\mathrm{coNP}$.
\end{theorem}
\begin{proof}
For any presented NP language $A$, take its supplied completeness reduction
$f$ to $K$. Equality $A(x)=K(f(x))$ gives the equality of complement bits.
The coNP presentation of $K$ means that $\bar K$ has a polynomial verifier.
Compose it with $f$ to obtain a verifier for $\bar A$. Thus complementation
maps every NP presentation to another NP presentation. A coNP presentation
is, by definition, the complement of an NP presentation, so the same
transformation gives both inclusions. If $\mathrm P=\mathrm{NP}$, take
$K=\mathrm{SAT}$: a polynomial decider and its output-bit reversal give P,
hence NP, presentations for both $K$ and $\bar K$.
\end{proof}

\begin{theorem}[N18: parsimonious completeness for witness counting]\label{thm:N18}
For the counting convention
\[
 f_V(x)=\sum_{y<2^{q(|x|)}} [V_T(x,y)=1],
\]
counting satisfying assignments of circuits and of 3CNF formulas is
\#$\mathrm P$-complete under parsimonious polynomial many-one reductions.
Counting indexed subsets of a positive integer list with a specified sum is
also \#$\mathrm P$-complete under those reductions.
\end{theorem}
\begin{proof}
N03 preserves the accepting witness bit strings exactly. N04 gives each such
string one and only one satisfying gate extension, and retains all input
variables. Therefore $f_V(x)$ equals the number of satisfying assignments of
$F_{C_x}$, with no multiplicative factor. Circuit counting has the same
identity before Tseitin conversion. Both counts belong to \#$\mathrm P$
because their finite assignments are verified in polynomial time. Applying
the bijection of N09 gives equality with the number of indexed subsets
summing to its target. This target count is in \#$\mathrm P$ by its binary
subset verifier. Bounded induction on the finite witness lists proves that
inverse bijections preserve cardinalities. This does not infer that exact
counting is polynomial time, and it does not identify arbitrary padded
nondeterministic path counts with the fixed-witness convention.
\end{proof}

\begin{theorem}[N19: a factor-two vertex-cover approximation]\label{thm:N19}
There is a polynomial algorithm that returns a vertex cover $C$ satisfying
$|C|\le2|C^*|$ for every vertex cover $C^*$ of the same graph.
\end{theorem}
\begin{proof}
Scan the edges in a fixed order. If both endpoints are unused, add the edge
to a matching $M$ and mark its endpoints. Output all marked vertices. An
uncovered edge would have had two still-unmarked endpoints when scanned,
because marks only accumulate. It would therefore have been selected, a
contradiction; hence the output is a cover. The selected edges have disjoint
endpoints, so $|C|=2|M|$. Every cover $C^*$ contains at least one endpoint of
each selected edge. Choose the smaller numbered endpoint belonging to $C^*$
on each such edge. Disjointness makes this choice injective, proving
$|M|\le|C^*|$. The algorithm never needs to compute an optimum cover; the
inequality holds for each supplied cover witness and hence for the optimum
obtained, if desired, by finite minimization. A single edge scan with simple
mark updates is polynomial on the explicit graph encoding.
\end{proof}

\begin{theorem}[N20: vertex cover is fixed-parameter tractable]\label{thm:N20}
For a graph on $N$ vertices and parameter $k$, vertex cover of size at most
$k$ can be decided, and a cover produced, in time
$2^{k+1}\poly(N+|E|+1)$.
\end{theorem}
\begin{proof}
If there are no edges, return the empty cover. If an edge remains and $k=0$,
return failure. Otherwise choose the first edge $\{u,v\}$. Every cover
contains $u$ or $v$. Recurse on the graph with $u$ deleted and parameter
$k-1$, and on the graph with $v$ deleted and parameter $k-1$. A successful
branch returns its chosen endpoint together with the recursively produced
cover. If both branches fail, a proposed cover for the original graph would,
by its choice of endpoint, restrict to a cover in a failed branch, a
contradiction. This proves correctness by induction on the remaining depth,
with all covers bounded by $2^N$ and all graphs by their fixed matrix bound.

The recursion tree has at most $1+2+\cdots+2^k<2^{k+1}$ nodes. At each node
one performs a polynomial scan and graph update. It is a primitive-recursive
algorithm: enumerate the full binary tree through depth $k$ and compute
bottom-up, or use an explicit stack with the stated node clock. For $k\ge N$
return all vertices immediately. The parameter $k$ is thereby separated from
the input-size polynomial, which is the fixed-parameter conclusion.
\end{proof}

\begin{theorem}[N21: UNSAT and tautology are coNP-complete]\label{thm:N21}
UNSAT, propositional TAUT, and tautology for DNF formulas with terms of width
at most three are coNP-complete under polynomial many-one reductions.
\end{theorem}
\begin{proof}
The complement of UNSAT is SAT. The complement of TAUT has an NP witness: an
assignment on which the formula is false. These are bounded Boolean
complements, so finite assignment enumeration also identifies them exactly,
without using double-negation elimination for an arbitrary proposition.
For any coNP presentation $A$, its complement has an NP reduction $f$ to 3SAT.
Hence $A(x)=\operatorname{UNSAT}(f(x))$. Negate that 3CNF syntactically and
apply De Morgan laws to obtain a DNF whose terms have width at most three.
For every assignment the resulting DNF has the opposite truth bit, so
$f(x)$ is unsatisfiable exactly when this DNF is a tautology. Negation and
De Morgan conversion have linear output size for a CNF. The bounded universal
assignment test follows from the corresponding pointwise equality by a
finite conjunction fold. This proves the remaining hardness claims.
\end{proof}

\begin{theorem}[N22: Horn SAT is P-complete under logspace reductions]\label{thm:N22}
Horn satisfiability is P-complete, even when every clause has width at most
three. The hardness reduction uses only finite gate-local rules.
\end{theorem}
\begin{proof}
Take an acyclic circuit with its input values specified, as in N06. For each
node $i$ introduce two variables $t_i,f_i$, intended to assert that its value
is respectively true or false. For an input or constant node, add the unit
fact for its specified value. For a NOT gate $z=\neg a$, add the rules
$f_a\to t_z$ and $t_a\to f_z$. For an AND gate $z=a\wedge b$, add
\[
 t_a\wedge t_b\to t_z,\qquad f_a\to f_z,\qquad f_b\to f_z.
\]
For an OR gate $z=a\vee b$, add
\[
 t_a\to t_z,\qquad t_b\to t_z,\qquad f_a\wedge f_b\to f_z.
\]
These are Horn clauses of width at most three. Finally add the negative unit
clause $\neg f_o$, where $o$ is the output node.

Gate-order induction proves that the least closure of the positive-head
rules contains exactly one of $t_i,f_i$ at every node, namely the assertion
of its actual evaluated bit. The base case is the supplied input facts.
For each gate, the displayed rules force its correct assertion from its
already determined input assertions, and none of them triggers its incorrect
assertion. No other rule has that node as head. Every model must contain
this least closure, by the same addition invariant used in N14. If the
circuit outputs zero, $f_o$ is forced and the final negative unit is
impossible. If it outputs one, the least closure itself satisfies every
rule and $\neg f_o$, giving a model. Thus the Horn instance is satisfiable
exactly when the circuit value is one.

A constructor scans the gate list and emits at most three clauses per gate,
plus the input facts and final unit. Variable indices are the two integers
$2i$ and $2i+1$, so its counters and address arithmetic need only logarithmic
work space. Compose with N06 to obtain hardness for P. N14 gives polynomial
membership, and the bounded gate-prefix invariant gives the object-arithmetic
correctness proof. The construction makes no unsupported claim that all
Horn models have exactly one assertion per node; only their least closure
has that property, and that is sufficient for the reduction.
\end{proof}

\subsection{The object-arithmetic obligations in this section}

Each constructor above is a fixed algorithm on finite codes, and each proof
uses finite lists with displayed bounds. To lower the proofs to the repository's
object theory, instantiate the Boolean bounded-formula reflection and bounded
induction interfaces; do not reinterpret a semantic Lean function parameter as
an object-language variable. The especially useful invariants are the
following.

\begin{itemize}
\item For circuit evaluation and Tseitin: every gate below the current prefix
has its computed value; all assignment codes are below $2^{q+G}$.
\item For machine unrolling: all time indices are at most $T$, all tape
positions below $L$, and all witnesses below $2^q$; the complete configuration
array agrees bit by bit with the bounded machine run.
\item For graph reductions: every selected occurrence lies in its recorded
clause; selected literals have consistent requested bits; vertex and subset
codes are bounded by their explicit vertex counts.
\item For subset sum: all digit sums are below eight, exactly one variable
number is selected, and the slack digit equals four minus the true-occurrence
count. Quotient/remainder induction gives the no-carry equality.
\item For 2SAT: the reachability table at stage $j$ matches paths with allowed
intermediate vertices below $j$; representatives, predecessor counts, and keys
are all bounded finite scans.
\item For Horn SAT and SAT search: the true-variable set is contained in each
model below $2^v$, or the current partial assignment has a satisfying extension
below $2^v$, respectively.
\item For parameterized vertex cover: every branch state has a graph below the
original adjacency bound and a chosen-vertex list of length at most $k$; all
branches lie in the full tree with fewer than $2^{k+1}$ nodes.
\end{itemize}

The accompanying \texttt{lean/ReductionsOutline.lean} supplies finite semantic
definitions and precise target statements. It is an uncompiled outline, not a
certificate for these theorems. The Python companion checks the new witness
maps and algorithms against exhaustive small reference computations, including
unused inputs, duplicated literals, empty clauses, and equal-valued indexed
subset entries. Such checks validate implementation cases; the general proofs
are the arguments above.
